% Módulo 01 — Introdução (R775+; problema, justificativa, estrutura)
\chapter{Introdução}\label{cap:intro}
\section{O problema}
Sistemas que delegam tarefas a agentes de software decidem sob incerteza com critérios múltiplos, e essas decisões raramente carregam garantias auditáveis: não se sabe demonstrar que os pesos somam um, que a regra escolhe o melhor pelos critérios declarados, nem que a medição de qualidade significa o que diz. A engenharia de software experimental estabeleceu há décadas que alegações exigem desenho, medição e relato disciplinados \cite{basili1986,wohlin2012,runeson2009}, e que revisões sistemáticas exigem protocolo de relato \cite{page2021}; a crise de reprodutibilidade documentada em múltiplas áreas \cite{osc2015,ioannidis2005,collberg2016} mostra o custo de evidência fraca. Este trabalho aplica esse padrão a um roteador por atenção: cada propriedade alegada ou é provada por máquina, ou é medida com incerteza declarada, ou é renunciada explicitamente.
\section{Questão e objetivos}
Que propriedades do núcleo matemático de um roteador inspirado em atenção podem ser formalizadas, e que comportamento sua implementação apresenta em uma avaliação sintética interna? Objetivo geral: formalizar e avaliar internamente um roteamento por atenção. Objetivos específicos: (OE1) verificar por máquina o núcleo matemático em Lean 4 com Mathlib \cite{moura2021,mathlib2020}; (OE2) comparar arrependimento e taxa de acerto com os comparadores, relatando intervalos exploratórios, seguindo a recomendação de estimação em vez de significância pontual \cite{cumming2014,wasserstein2016}; (OE3) auditar axiomas, alegações, referências e código, incluindo inspeção e métricas na tradição de \cite{fagan1976,mccabe1976,chidamber1994}; (OE4) declarar limites e publicar a trilha de reprodução segundo princípios FAIR \cite{wilkinson2016} e documentação de dados e modelos \cite{gebru2021,mitchell2019}.
\section{Justificativa}
Roteamento opaco não é auditável; garantias formais exigem hipóteses explícitas; medições requerem análise de incerteza. Como perspectiva educacional, trilhas auditáveis desse tipo oferecem material didático para o ensino de avaliação e validação --- hipótese de trabalho futuro, não resultado desta dissertação.
\section{Delimitação e tese defendida}
Defende-se exclusivamente otimalidade relativa à utilidade definida e evidência preliminar interna. Não se compara com plataformas existentes, não se generaliza além do simulador descrito e não se transfere propriedades de Transformers treinados. A tese metodológica é que verificação por máquina, teste numérico de propriedades, ablação e auditoria de alegações compõem um padrão de evidência auditável mesmo sem validação externa.
\section{Estrutura do trabalho}
A Seção~\ref{cap:preliminares} fixa preliminares matemáticos; a Seção~\ref{cap:fundamentacao} fundamenta atenção, correspondência Q/K/V, combinação convexa, seleção gulosa, Lean e os trabalhos relacionados; a Seção~\ref{cap:metodo} detalha método, estatística, ameaças e validade; a Seção~\ref{cap:provas} apresenta as onze provas comentadas; a Seção~\ref{cap:resultados}, os testes de propriedades e a avaliação empírica; a Seção~\ref{cap:discussao} discute revisões e riscos; a Seção~\ref{cap:arguicoes} antecipa arguições e responde nos limites do texto; a Seção~\ref{cap:conclusao} conclui; o glossário define os termos principais; os apêndices preservam código, resultados, transcrições e especificações.
\section{Cenário-guia de roteamento (saídas registradas)}
Tarefa ``provar teorema formal com código Lean'' exigindo \texttt{formal} e \texttt{codigo}, com três agentes: Ana (prova e código; carga 0,1; confiança 0,9), Bia (código e dados; carga 0,8; confiança 0,5) e Cid (as três capacidades; carga 0,4; confiança 0,7). O gate exclui Bia (falta \texttt{formal}); utilidades: Ana $0{,}9175$, Cid $0{,}8161$; pesos finais Ana $0{,}5253$ e Cid $0{,}4747$ (soma 1). Ana vence com menos capacidades porque combina menor carga, maior confiança e melhor ajuste semântico --- exatamente o tipo de compensação que os teoremas disciplinam e a bancada mede.
\par Dois casos-limite reais: (i) agentes gêmeos que diferem só na descrição recebem $0{,}5053$ contra $0{,}4947$ --- nem gêmeos empatam na semântica, e empates exatos decidem-se pelo identificador, deterministicamente; (ii) candidato único recebe peso $1{,}0$ --- o T7 em execução.
